\begin{trapbox}
	\begin{description}[style=nextline, leftmargin=2.8em, labelsep=0.5em, itemsep=0.35em]
		\item[\textbf{\textcolor{CTrap}{易错点一（忽视定义域）}}] 错误地认为对于任意正实数 $a,b$，不等式 $\dfrac{\ln a}{\ln b}<\dfrac{\ln(a+c)}{\ln(b+c)}$ 都成立。
		\item[\textbf{\textcolor{CTrap}{正确理解（要求a>1,b>1）：}}] 只有当 $a>1$ 且 $b>1$ 时，$\ln a,\ln b$ 为正，不等式链才成立。
		\item[\textbf{\textcolor{CTrap}{错因分析（符号忽视）：}}] 当 $0<a<b<1$ 时，$\ln a,\ln b$ 为负，此时原不等式不再成立，不能直接套用。
	\end{description}

	\medskip
	\begin{description}[style=nextline, leftmargin=2.8em, labelsep=0.5em, itemsep=0.35em]
		\item[\textbf{\textcolor{CTrap}{易错点二（增量c的符号）}}] 错误地将 $c$ 视为任意实数，认为只要加上同一个数即可。
		\item[\textbf{\textcolor{CTrap}{正确理解（要求c>0）：}}] 必须 $c>0$，才能保证 $a+c>a$ 且函数 $H(x)=\dfrac{\ln x}{\ln(x+c)}$ 单调递增。
		\item[\textbf{\textcolor{CTrap}{错因分析（条件遗漏）：}}] 若 $c<0$，则 $a+c<a$，不等式方向可能改变，原结论不再成立。
	\end{description}

	\medskip
	\begin{description}[style=nextline, leftmargin=2.8em, labelsep=0.5em, itemsep=0.35em]
		\item[\textbf{\textcolor{CTrap}{易错点三（取倒数方向）}}] 在利用取倒数变形时，忘记将不等号反向。
		\item[\textbf{\textcolor{CTrap}{正确理解（方向反转）：}}] 当不等式两边均为正时，取倒数后不等号严格反向，例如 $\frac{\ln3}{\ln4}<\frac{\ln4}{\ln5}$ 推出 $\frac{\ln4}{\ln3}>\frac{\ln5}{\ln4}$。
		\item[\textbf{\textcolor{CTrap}{错因分析（法则混淆）：}}] 部分学生习惯乘法除法，对取倒数改变不等号方向不熟练。
	\end{description}
\end{trapbox}
